\section{Tám chỉ số trước ground truth}
\label{sec:ch09-ket-qua}

\index{AUROC}%
\index{chỉ số trung thực!kết quả kiểm định}%
Tám chỉ số được đem ra chấm, và chúng chính là bốn nhóm mà
mục~\ref{sec:ch08-dieu-chua-giai-quyet} đã mượn cách chia. Nhóm dựa trên mức
quan trọng gồm Adding Mistakes, Early Answering, Filler Tokens và SCM; nhóm dựa
trên \tn{tính hữu dụng ngữ nghĩa}{semantic utility} gồm Paraphrasing và
Simulatability; nhóm dựa trên
tham số chỉ có FUR; nhóm dựa trên attribution chỉ có
CC-SHAP~\autocite{p21metrics}. Cần nói ngay một điều về hình thức trình bày:
bài báo không viết công thức cho bất kỳ chỉ số nào trong tám chỉ số ấy mà mô tả
tất cả bằng lời, nên sách cũng không phát biểu công thức nào thay nó.

Hai dòng đối chứng được đưa vào để đọc bảng, và chúng không cùng vai. Random
gán nhãn theo lối bốc thăm. LM Judge (generic) là một mô hình ngôn ngữ được
giao vai người giám sát chuỗi suy luận nói chung, không kèm định nghĩa của bài
báo. Dòng LM Judge còn lại thì được cấp đúng hai định nghĩa của
mục~\ref{sec:ch09-dinh-nghia-lai}, và các tác giả xếp nó là trần trên chứ không
phải một chỉ số dự thi, vì bộ dữ liệu này được dựng bằng chính các đường ống
dùng mô hình ngôn ngữ chạy trên cùng dữ liệu ấy, nên một giám khảo cùng loại
được hưởng đúng nguồn tín hiệu đã sinh ra nhãn~\autocite{p21metrics}. Đọc con
số của nó thành bằng chứng rằng mô hình ngôn ngữ giải xong bài toán này là đọc
ngược lại bài báo.

\begin{table}[htbp]
  \centering
  \footnotesize
  \caption{AUROC kèm biên 95\% theo DeLong, chép từ bảng kết quả của bài
  báo~\autocite{p21metrics}. Cột \enquote{bước} chấm từng bước một, cột
  \enquote{chuỗi} chấm cả chuỗi; \enquote{không chạy} nghĩa là bài báo không
  chạy biến thể ấy ở mức đó, chứ không phải chạy rồi bỏ kết quả. Hai dòng cuối
  là đối chứng, và LM Judge là trần trên chứ không phải chỉ số dự thi.}
  \label{tab:ch09-auroc}
  \begin{tabular}{@{}llll@{}}
    \toprule
    Nhóm & Chỉ số & Bước & Chuỗi \\
    \midrule
    Mức quan trọng & Adding Mistakes & $0.51 \pm 0.02$ & $0.51 \pm 0.04$ \\
                   & Early Answering & $0.51 \pm 0.01$ & $0.45 \pm 0.03$ \\
                   & Filler Tokens   & $0.59 \pm 0.01$ & $0.50 \pm 0.02$ \\
                   & SCM             & không chạy      & $0.38 \pm 0.03$ \\
    \addlinespace
    Hữu dụng ngữ nghĩa & Paraphrasing    & không chạy & $0.61 \pm 0.03$ \\
                       & Simulatability  & không chạy & $0.50 \pm 0.01$ \\
    \addlinespace
    Tham số     & FUR     & $0.52 \pm 0.02$ & không chạy \\
    Attribution & CC-SHAP & $0.41 \pm 0.03$ & $0.70 \pm 0.04$ \\
    \midrule
    Trần trên & LM Judge           & $0.87 \pm 0.01$ & $0.82 \pm 0.02$ \\
    Đối chứng & Random             & $0.5 \pm 0$     & $0.5 \pm 0$ \\
              & LM Judge (generic) & $0.68 \pm 0.02$ & $0.67 \pm 0.04$ \\
    \bottomrule
  \end{tabular}
\end{table}

Bảng~\ref{tab:ch09-auroc} là kết quả trung tâm của cả Phần~IV, và cách đọc nó
chỉ gồm một phép so: cột nào cũng phải đặt cạnh 0.5. Phần lớn các ô nằm
sát con số ấy. Chỉ số cao nhất ở mức cả chuỗi là CC-SHAP với 0.70, còn ở mức
từng bước là Filler Tokens với 0.59~\autocite{p21metrics}. Đó là hai thành
tích tốt nhất mà cả một lĩnh vực đề xuất chỉ số đạt được khi lần đầu bị đem ra
đối chiếu với ground truth.

Hai thành tích ấy còn không giữ được qua hai mức đo. CC-SHAP tụt xuống 0.41 ở
mức từng bước, tức là thấp hơn bốc thăm, trong khi Filler Tokens về sát mức
bốc thăm ở mức cả chuỗi~\autocite{p21metrics}. Một dụng cụ đo tốt ở thang này
và tệ hơn bốc thăm ở thang kia thì chưa đo được cái nó ghi trên mặt đồng hồ;
nó đang phản ứng với một thứ khác mà ở một trong hai thang thì thứ đó tình cờ
tương quan với nhãn.

\index{độ chệch dự đoán}%
Các con số sát 0.5 ấy không phải là kiểu \enquote{đoán hú họa cân bằng}. Trên
một mẫu đã cân bằng nhãn, nhóm dựa trên mức quan trọng gán nhãn không trung
thực cho từ 90\% tới 96\% số trường hợp ở mức cả chuỗi, còn nhóm dựa trên tính
hữu dụng ngữ nghĩa gán nhãn trung thực cho từ 94\% tới
96\%~\autocite{p21metrics}. Nói cách khác, mỗi nhóm gần như luôn trả lời một
đáp án cố định, và AUROC quanh 0.5 là hệ quả của
\tn{độ chệch dự đoán}{prediction skew} ấy chứ không phải của nhiễu.

Nếu các chỉ số cùng đo một đại lượng thì dù có sai chúng cũng phải sai giống
nhau, nên mức đồng thuận giữa chúng là một phép kiểm độc lập. Đo bằng kappa
của Cohen, phần lớn các cặp nằm quanh 0; ở mức cả chuỗi cặp cao nhất đạt 0.35,
còn ở mức từng bước cặp cao nhất chỉ đạt 0.12. Các tác giả đọc nó đúng như
vậy: mức đồng thuận trần ở 0.35 cho thấy chúng đo những tính chất khác
nhau~\autocite{p21metrics}.

Còn một trục nữa mà bảng~\ref{tab:ch09-auroc} không thể hiện, và với người
định dùng các chỉ số này để giám sát thì nó quyết định. CC-SHAP, tức chỉ số
tốt nhất ở mức cả chuỗi, tốn tới $10^3$ giây cho mỗi trường hợp. Biến thể mức
cả chuỗi của FUR thì không được chạy lần nào, vì các tác giả ước tính nó cần
tới $10^5$ giây cho mỗi chuỗi, tức hơn một ngày~\autocite{p21metrics}. Chỉ số
đắt nhất trong bảng lại là chỉ số duy nhất nhỉnh hơn hẳn mức bốc thăm, và một
ô của bảng bỏ trống chỉ vì không ai trả nổi chi phí đo nó.

Bảng cho biết các chỉ số hỏng. Vì sao chúng hỏng theo đúng kiểu chúng hỏng thì
là việc của mục sau.
